\documentstyle[%


righttag%


,11pt%


,amssymb%


,russian%               remove in Eng.


,izvru%                 Set izven% in Eng.


]{amsart}





%       Math definitions


% NB: repl. ... by \dotc in Eng.


\newcommand{\col}{\operatorname{col}}


\newcommand{\Def}{\overset{\mathrm{def}}{=}{}}


\newcommand{\tts}[1]{}





%\hoffset=-15mm%                  For Eng.


\setcounter{page}{3}


\god{1999}


\nomer{10~(449)} %                        Remove in Eng.


%\nomer{Vol.\,43, No.\,10}%               Set in Eng.


%\str{\pageref{begin}--\pageref{end}}%  Set in Eng.


\udk{517.929.4\,:\,519.21}





\author{\it Р.И.\,Кадиев}


% NB:   ^^ remove in Eng.


\title{К вопросу об устойчивости стохастических


функционально-дифференциальных уравнений\\


по первому приближению}





\date{}


%\pagestyle{myheadings}%         Set in Eng.


\pagestyle{plain} %              Remove in Eng.


\begin{document}


%\thispagestyle{plain}%          Set in Eng.


\maketitle


\label{begin}








Работа посвящена распространению известной


теоремы Ляпунова об устойчивости по первому приближению на


случай функционально-дифференциальных уравнений


по семимартингалу. Для детерминированных


функционально-дифференциальных уравнений этот


вопрос изучался в [1].





Пусть $(\Omega, {\cal F}, ({\cal F})_{t\ge0},P)$ ---


стохастический базис ([2], гл.\,1, \S\,1, c.\,9);


$D^n$ --- линейное пространство $n$-мерных


предсказуемых ([2], гл.\,1, \S\,1, c.\,13) случайных процессов


на $[0, \infty [$, траектории которых почти наверно


(п.н.) непрерывны справа и имеют предел слева;


$k^n$ --- линейное пространство $n$-мерных


${\cal F}_0$-измеримых случайных величин;


$Z=\col(z_1,\dotsc,z_m)$ есть $m$-мерный семимартингал


([2], гл.\,1, \S\,1, c.\,73); $L(Z)$ --- линейное пространство


предсказуемых $n\times m$-матриц на $[0, \infty [$, строки


которых локально интегрируемы по семимартингалу


$Z$ [3]; $\lambda : [0, \infty [ \rightarrow R_+ $ --- некоторая


неубывающая функция; $\lambda ^*$ --- мера,


порожденная функцией $\lambda $; $L^\lambda _q$ ---


линейное пространство скалярных функций на $[0, \infty [$,


суммируемых со степенью $q$ ($1 \le q < \infty$) по функции


$\lambda$ и ограниченных в существенном по мере


$\lambda ^*$ при $q = +\infty $ $(L^t_q = L_q)$;


$R^n$ --- линейное пространство $n$-мерных векторов


с нормой $|\cdot|$; $\|\cdot\|$ --- норма $k\times n$-матрицы,


согласованная с нормой $|\cdot|$; $\|\cdot\|_X$ --- норма


в нормированном пространстве $X$; $1 \le p < \infty$;


$E$ --- символ математического ожидания.





Для заданного числа $p$ и положительной скалярной


функции $\gamma (t)$ $(t \in [0, \infty [)$ введем


следующие линейные пространства:


\begin{gather*}


M^\gamma _p = \{x: x \in D^n,\ \|x\|_{M^\gamma _p}\Def


(\sup_{t\ge0}E|\gamma (t)x(t)|^p)^{1/p} < \infty \};\\


%


k^n _p = \{\alpha : \alpha \in k^n,\ \|\alpha \|_{k^n _p}\Def


(E|\alpha |^p)^{1/p} < \infty \}


\end{gather*}


и $M^1_p = M_p$.





Если $B$ --- линейное подпространство


пространства $L(Z)$ c нормой $\|\cdot\|_B$, то положим


$B^\gamma = \{f: f \in B,\ \gamma f \in B\}$.


Ясно, что $B^\gamma $ --- пространство с нормой


$\|f\|_{B^\gamma } = \|\gamma f\|_B$.





Рассмотрим уравнение


\begin{equation}


dx(t) = [(Qx)(t) + g(t)]dZ(t),\quad t \ge 0,


\end{equation}


где $g \in L(Z)$, а $Q:D^n \to L(Z )$ --- $k^1$-линейный


[4] оператор. Уравнение (1) рассматривается в предположении,


что для любого $\alpha \in k^n$ существует единственное


(с точностью до $P$-эквивалентности)


решение $x(t)$ c $x(0) = \alpha$.


Для этого решения имеет место представление


\begin{equation}


x(t) = U(t)x(0) + (Wg)(t),\quad t \ge 0,


\end{equation}


где $U(t)$ --- фундаментальная матрица, а $W$ --- оператор


Коши ($k^1$-линейный оператор, действующий из


пространства $L(Z)$ в пространство $D^n$) для


уравнения (1) [4].





Обозначим через $D^n_p$ линейное пространство,


определенное равенством (2) при $g \in B$, $x(0) \in k_p^n$.


Здесь и в дальнейшем $B$ --- линейное нормированное


подпространство пространства $L(Z)$. Таким образом,


$D^n_p = WB\oplus Uk^n_p$. Если $D_p^n \subset M_p^\gamma$


при некоторой функции $\gamma$,


то, очевидно, $D_p^n$ --- линейное нормированное


пространство. В дальнейшем будем предполагать, что


$D_p^n \subset M_p^\gamma $ при некоторой функции $\gamma$


и $\wh D_p^n$ ---


замыкание пространства $D_p^n$.





Пусть $V:D_p^n \to B$ --- $k^1$-линейный


оператор. Если задача Коши


\begin{gather}


dx(t) = [(Vx)(t) + f(t)]dZ(t),\quad t \ge 0,\\


%


x(0) = \alpha


\end{gather}


имеет единственное решение (с точностью до


$P$-эквивалентности) $x \in D_p^n$ при любых


$f \in B$ и $\alpha \in k_p^n$, то существуют такие


$k^1$-линейные операторы $K:B \to D_p^n$,


$X: k_p^n \to D_p^n$, что для решения


задачи (3), (4) имеет место представление


\begin{equation}


x(t) = (Kf)(t) + X(t)\alpha,\quad t \ge 0.


\end{equation}





Обозначим через $x_f(t,\alpha )$ решение задачи


(3), (4), а через $\Theta\Def


W(V - Q)$ --- $k^1$-линейный оператор, действующий


из пространства $D^n_p$ в пространство $D^n_p$.


Тогда в силу эквивалентности задачи (3), (4) и уравнения


$x(t)=U(t)\alpha+(\Theta x)(t)+(W f)(t)$, $t\ge 0$, 


справедлива





\begin{lem} Задача Коши $(3)$, $(4)$ имеет единственное


решение $x \in D^n_p$ для каждой пары


$\{f,\alpha \} \in B\times k_p^n$ тогда и только


тогда, когда существует обратный


оператор для оператора $(I - \Theta ): D_p^n\to D_p^n$.


\end{lem}





Отметим, что если $\|\Theta \|_{D_p^n\to D_p^n} < 1$


и оператор $\Theta $ действует из пространства $\wh D_p^n$


в пространство $D_p^n$, то оператор $(I - \Theta ): D_p^n


\to D_p^n$ обратим.





\begin{defn}


Будем говорить, что


уравнение (3) $D_p^n$-устойчиво, если задача


(3), (4) имеет единственное решение $x_f(\cdot,\alpha)\in D_p^n$


при любых $f \in B$, $\alpha \in k_p^n$,


и это решение непрерывно зависит от $\{f,\alpha \}$,


т.е. для любого $\varepsilon > 0$ существует такое $\delta > 0$,


что $\|x_{f_1}(\cdot,\alpha _1)


- x_f(\cdot,\alpha _2)\|_{D_p^n} < \varepsilon$,


если $\|f_1 - f\|_B < \delta$, $\|\alpha_1-\alpha\|_{k_p^n}<\delta$.


\end{defn}





Таким образом, $D_p^n$-устойчивость уравнения (3) --- это


корректная разрешимость задачи (3), (4).





\begin{lem}


Если уравнение $(3)$ $D_p^n$-устойчиво,


то для этого уравнения допустима пара $(M_p^\gamma ,B)$ $[4]$.


\end{lem}





\begin{pf}


Пусть уравнение (3) $D_p^n$-устойчиво.


Необходимо доказать допустимость пары $(M_p^\gamma ,B)$


для уравнения (3). Предположим противное, т.е. для уравнения (3)


пара $(M_p^\gamma ,B)$ не является допустимой. Тогда хотя бы


один из операторов $X:k_p^n \to D_p^n$ или


$K:B \to D_p^n$ для уравнения (3) неограничен.


Отсюда следует, что уравнение (3) не является


$D_p^n$--устойчивым. Следовательно, предположение неверно,


т.е. для уравнения (3) допустима пара $(M_p^\gamma ,B)$.


\end{pf}





Обозначим через $D_p^n(V)$ линейное пространство


решений задачи (3)--(4), определяемое равенством (5)


при $f \in B$, $\alpha \in k_p^n$. Для $D_p^n$-устойчивого


уравнения (3) имеет место $D_p^n = D_p^n(V)$. Это


равенство характеризует определенные особенности


асимптотического поведения решений уравнения (3) для начальных


условий $x(0) \in k_p^n$ и


$f \in B$. Отметим, что из $D_p^n$-устойчивости


уравнения (3) следует также и $M_p^\gamma $-устойчивость


[5] этого уравнения.





Распространяя сказанное на нелинейное уравнение


\begin{equation}


dx(t) = [(Nx)(t) + f(t)]dZ(t),\quad t \ge 0,


\end{equation}


с оператором $N: D_p^n \to B$ ($f \in B$), можно определить


$D_p^n$-устойчивость уравнения (6) как корректную


разрешимость в пространстве $D_p^n$ задачи Коши для этого уравнения с


начальными условиями


\begin{equation}


x(0) = \alpha,


\end{equation}


где $\alpha \in k_p^n$. Однако такая устойчивость


``в целом'' --- свойство исключительное для нелинейного уравнения и


требует от оператора $N$ условий, не выполнимых даже в простых


случаях. В детерминированном случае это подтверждается, в частности,


примером из [1]. Поэтому, как и в


детерминированном случае, остановимся на вопросе об


``устойчивости тривиального решения'' уравнения, к которому,


как известно, сводятся более общие случаи.





Аналогично случаю задачи (3), (4) через $x_f(t,\alpha )$


обозначим решение задачи (6), (7).





\begin{defn}


Пусть $(N0)(t) = 0$. Уравнение


(6) назовем локально


$D_p^n$-устойчивым, если можно указать


такое $\delta _0 > 0$, что для каждой пары $\{f,\alpha \} \in B\times


k_p^n$,


удовлетворяющей условиям $\|f\|_B < \delta _0$, $\|\alpha \|_{k_p^n} <


\delta _0$, задача (6)--(7) имеет единственное решение


$x_f(\cdot,\alpha ) \in D_p^n$ и для любого $\varepsilon > 0$ существует


такое $\delta = \delta (x_f(\cdot,\alpha ),\varepsilon) > 0$, что


$\|x_{f_1}(\cdot,\alpha ) - x_f(\cdot,\alpha )\|_{D_p^n} < \varepsilon$,


если $\|f_1 - f\|_B < \delta$, $\|\alpha _1 - \alpha \|_{k_p^n} < \delta$


и $\|f_1\|_B < \delta _0$, $\|\alpha _1\|_{k_p^n} < \delta _0$.


\end{defn}





Отметим, что если уравнение (6) локально $D_p^n$-устойчиво,


то тривиальное решение уравнения $dx(t) = (Nx)(t)dt$, $t \ge 0$,


$p$-устойчиво, если $\gamma (t) = 1$, асимптотически


$p$-устойчиво, если $\underset{t\to\infty}{\varliminf}{}


\gamma (t) = \infty$, и экспоненциально


$p$-устойчиво, если $\gamma (t) = \exp \{\beta t\}$, $\beta > 0$.





Пусть $Nx = Vx + Fx$, где $V:D_p^n \to B$ ---


$k^1$-линейный оператор, а нелинейный оператор


$F: \wh D_p^n \to B$ обладает свойством $(F0)(t) = 0$.


Тогда имеет место





\begin{thm}


Пусть уравнение $(3)$


$D_p^n$-устойчиво и для любого $k > 0$ найдется такое


$\delta > 0$, что


\begin{equation}


\|Fx_2 - Fx_1\|_B \le k\|x_2 - x_1\|_{D_p^n}


\end{equation}


при всех $x_1, x_2 \in D_p^n$, $\|x_1\|_{D_p^n} \le \delta$,


$\|x_2\|_{D_p^n} \le \delta$.


Тогда уравнение $(6)$ локально $D_p^n$-устойчиво.


\end{thm}





\begin{pf}


Пусть $X:k_p^n \to D_p^n$,


$K:B \to D_p^n$ --- операторы, определяющие


представление (5) решения задачи (3)--(4),


$\wh b\Def\|X\|_{k_p^n\to D_p^n}$,


$\wh c\Def \|K\|_{B\to D_P^n}$.


Задача Коши (6)--(7)


эквивалентна уравнению


\begin{equation}


x(t) = (KFx)(t) + h(t),\quad t \ge 0,


\end{equation}


в пространстве $D_p^n$, если $h(t) = (Kf)(t) + X(t)\alpha $.





В силу условий теоремы для некоторого положительного $k < 1/\wh c$


найдется такое $\varepsilon _0 > 0$, что при любых $x_i \in D_p^n$,


$\|x_i\|_{D_p^n} \le \varepsilon _0$, $i = 1,2$, имеет место


неравенство (8). Таким образом, в шаре $\{\|x\|_{D_p^n} \le \varepsilon _0\}$


при условии $\|h\|_{D_p^n} < \varepsilon _0(1 - \wh ck)$


ввиду локального принципа Банаха о сжимающих


отображениях ([6], 4.3.5) существует единственное решение


$x \in \wh D_p^n$ уравнения (9). В силу $D_p^n$-устойчивости


уравнения (3) и действия


оператора $F$ из пространства $\wh D_p^n$ в пространство


$B$ получим $x \in D_p^n$. Следовательно, задача (6)--(7)


имеет единственное решение из пространства $D_p^n$


при выполнении предыдущих условий. Кроме того, для


решений $x_{f_i}(\cdot,\alpha )$ этой задачи при условии,


что $\|f_i\|_B < \delta _0$, $\|\alpha _i\|_{k_p^n} < \delta _0$,


$i = 1,2$, $\delta _0 = \varepsilon _0(1 - \wh ck)/(\wh b + \wh c)$,


имеем


$$


\|x_{f_2}(\cdot,\alpha ) - x_{f_1}(\cdot,\alpha )\|_{D_p^n} \le


\|h_2 - h_1\|_{D_p^n}/(1- \wh ck) \le


(\wh c\|f_2 - f_1\|_B + \wh b\|\alpha _2 -


\alpha _1\|_{k_p^n})/(1 - \wh ck),


$$


где $h_i = Kf_i + X\alpha _i$, $i = 1,2$. Отсюда следует


непрерывная зависимость (по норме пространства $D_p^n$)


решения $x$ уравнения (9) от $h$ при


$\|h\|_{D_p^n} < \varepsilon _0(1- \wh ck)$,


а также задачи (6)--(7) от $f$ и $\alpha $ при $\|f\|_B < \delta _0$,


$\|\alpha \|_{k_p^n} < \delta _0$.


\end{pf}





Из доказательства теоремы 1 видно, что возможна также


следующая ее редакция.


\medskip





{\bf Теорема $\bold 1^{\boldsymbol\prime}$.} {\it Пусть уравнение


$(3)$ $D_p^n$-устойчиво,


$\wh c\Def \|K\|_{B\to D_p^n}$ ---


норма оператора Коши $K:{B \to D_p^n}$ этого уравнения,


$F:\wh D_p^n \to B$, $(F0)(t) = 0$ и существуют такие


постоянные $\delta > 0$, $0 < k < 1/\wh c$, что $\|Fx_2 - Fx_1\|_B \le


k \|x_2 - x_1\|_{D_p^n}$ для всех $x_i \in D_p^n$,


$\|x_i\|_{D_p^n} \le \delta$,


$i = 1,2$. Тогда уравнение $(6)$ локально $D_p^n$-устойчиво.}


\medskip





Очевидно, не для любых уравнений вида (1) и подпространства


$B$ пространства $L(Z)$ будет выполняться условие


$D_p^n \subset M_p^\gamma $


при некоторой функции $\gamma $. Рассмотрим некоторые уравнения


вида (1) и подпространства $B$ пространства $L(Z)$, для которых


$D_p^n \subset M_p^\gamma $ при некоторой функции $\gamma$.


Для этого в дальнейшем предположим,


что семимартингал $Z$ представим в виде $Z = b + c$, где


$b$ --- предсказуемый случайный процесс локально ограниченной


вариации, а $c$ --- локально квадратично интегрируемый


мартингал ([1], гл.\,1, \S\,1, с.\,28). Кроме того,


все компоненты процесса $b$ и взаимные характеристики


$\langle c^i,c^j\rangle$ ([2], гл.\,1, \S\,8, с.\,48) всех компонент


мартингала $c$ будем предполагать абсолютно непрерывными


относительно меры, задаваемой функцией


$\lambda$. Последнее означает, что


$$


b^i = \int_0^\cdot a^id\lambda, \quad


\langle c^i,c^j\rangle = \int_0^\cdot A^{i,j}d\lambda


\quad (i, j = 1,\dotsc,m).


$$


Пусть $a=\col(a^1,\dotsc,a^m)$, $A = [A^{ij}]$ --- $n\times m$-матрица.


Известно [3], что в этом случае пространство $L(Z)$ состоит из


предсказуемых $n\times m$-матриц $H$, для которых выполнено


неравенство $\int \limits _0^t(|Ha| +\|HAH^T\|)d\lambda < \infty$


п.н. для любого $t\ge0$ и $\int\limits_0^t H\,dZ=\int\limits_0^tH\,db


+ \int \limits _0^tH\,dc$. Кроме того, имеет место неравенство


$(E|\int \limits _0^tH\,dZ|^{2p})^{1/2p} \le (E(\int \limits


_0^t|Ha|d\lambda )^


{2p})^{1/2p} + c_p(E(\int \limits _0^t\|HAH^T\|d\lambda )^p)^{1/2p}$,


где $c_p$ --- положительное число, зависящее от $p$ ([2], гл.\,1, \S\,9,


с.\,65).





Рассмотрим уравнение вида


\begin{equation}


dx(t) + \ov \alpha \xi (t)x(t)d\lambda (t) = g(t)dZ(t),\quad t \ge 0,


\end{equation}


где $\xi $ --- скалярная неотрицательная функция на $[0, +\infty[$,


локально суммируемая по мере, задаваемой функцией $\lambda $, $\ov \alpha$


--- некоторое положительное число. Разумеется, выбор подпространств $B$


пространства $L(Z)$, для которых


$D_p^n \subset M_p^\gamma$ при некоторой функции $\gamma$,


существенно зависит от


$\xi$ и $\ov \alpha$. Рассмотрим примеры таких


подпространств, являющихся


разновидностью функциональных пространств со смешанной


нормой. Пусть $1 \le q \le \infty$, обозначим через


$$


\Lambda _{p,q}^n(\xi )=\{H: H \in L(Z),\ (E|Ha|^p)^{1/p}\xi ^{1/q-1},


\ (E\|HAH^T\|^{p/2})^{1/p}\xi ^{1/q-1/2} \in L_q^\lambda \}


$$


линейное подпространство $L(Z)$ с нормой


$$


\|H\|_{\Lambda^n_{p,q}(\xi)}\Def


\|(E|Ha|^p)^{1/p}\xi ^{1/q-1}\|_{L_q^\lambda} +


\|(E\|HAH^T\|^{p/2})^{1/p}\xi ^{1/q-1/2}\|_{L_q^\lambda }


$$


$(\Lambda _{p,q}^n(1) = \Lambda _{p,q}^n)$.


Тогда из теоремы 3 [4] следует, что если


$B = {(\Lambda _{2p,q}^n(\xi ))}^\gamma$,


$\gamma (t) = \exp\big\{\beta \int \limits _0^t\xi (s)ds\big\}$, $\beta>0$,


то $D_{2p}^n \subset M_{2p}^\gamma$.





Применим вышеизложенное к уравнению Ито с ``максимумом''.


В этом случае $Z(t) = \col(t,{\cal W}^1(t),\dotsc,{\cal W}^{(m-1)}(t))$,


где ${\cal W}^i$ ($i = 1,\dotsc,m-1$) --- независимые винеровские


процессы и $\lambda (t) = t$, $a = \col(1,0,\dotsc,0)$,


$A=\text{diag\,}[0,1,\dotsc,1]$.





Рассмотрим уравнение вида


\begin{equation}


dx(t) = [(Vx)(t) + (Fx)(t) + f(t)]dZ(t),\quad t \ge 0,


\end{equation}


где $f \in B$, $V:D_p^n \to B$ --- $k^1$-линейный оператор,


$Fx\Def T\{x,S_1^\varphi x,\dotsc,S_r^\varphi x\}$ и


$F:\wh D^n_p \to B$ --- нелинейный оператор,


а оператор $S_d^\varphi $ определен равенством


$$


S_d^\varphi x\Def


\col(S_d^{\varphi _1}x_1,\dotsc,S_d^{\varphi _n}x_n),


$$


где $(S_d^{\varphi_i}x_i)(t)\Def


\sup \limits _{s\in [h_{id}(t),g_{id}(t)]}{(E|\wh x_i(s)|^{2p})}^{1/2p}$,


$$


\wh x_i(s)\Def


\begin{cases}


x_i(s), &\text{если \;} s \ge 0;\\


\varphi _i(s), &\text{если \;} s < 0,


\end{cases}


$$


$x=\col(x_1,\dotsc,x_n)$, $\varphi=\col(\varphi_1,\dotsc,\varphi _n)$,


$h_{id}, g_{id}: [0, + \infty [ \to R^1$ --- измеримые


по Лебегу функции такие, что $h_{id}(t) \le g_{id}(t)$


при всех $t \ge 0$, $\varphi _i:\, ]- \infty , 0[


\times \Omega \to R^1$ ---


случайный процесс с п.н. непрерывными траекториями,


$\sup \limits _{t<0}{(E|\varphi _i(t)|^{2p})}^{1/2p} < \infty $ при


$i = 1,\dotsc,n$, $d = 1,\dotsc,r$.





В дальнейшем в качестве уравнения (1) возьмем


уравнение (10) с $\xi(t) \equiv t$,


$t \ge 0$, $B = (\Lambda _{2p,\infty }^n)^\gamma$,


$\gamma (t) = \exp \{\beta t\}$, $\beta \in [0, \ov \alpha [$,


$B^1 = {(L^n_\infty )}^\gamma = \{f: f \in L^n_\infty,


\ \gamma f \in L^n_\infty \}$ --- линейное пространство с нормой


$\|f\|_{B^1} = \|\gamma f\|_{L^n_\infty }$.


Тогда по теореме 3 из [4] $D_{2p}^n \subset M_{2p}^\gamma$.


\medskip





{\bf Замечание.} Так как рассматривается уравнение Ито, то в


дальнейшем все случайные процессы


прогрессивно измеримы относительно потока $\sigma $-алгебр


$({\cal F}_t)_{t\ge 0}$ ([2], гл.\,1, \S\,1, с.\,11).


Ясно, что любой


случайный процесс из пространства $\wh D_{2p}^n$ прогрессивно


измерим и его траектории п.н. непрерывны.





\begin{lem}


Для каждого случайного процесса $x \in \wh D_{2p}^n$


и для каждого $d$ $(1 \le d \le r)$ функция $S_d^\varphi x$


измерима.


\end{lem}





Доказательство леммы дословно


повторяет доказательство леммы 1 из [1]. При


этом необходимо воспользоваться тем, что функция


${(E|\wh x|^{2p})}^{1/2p}$ непрерывна справа, в чем


нетрудно убедиться непосредственной


проверкой.





Обозначим $S_d^0\Def


S_d^\varphi $ при $\varphi = 0$, $d = 1,\dotsc,r$.





\begin{lem}


Пусть существует такое $\delta > 0$,


что $0 \le t - h_{id}(t) \le 


\delta$ при всех $t \ge 0$ и всех


$i = 1,\dotsc,n$, $d = 1,\dotsc,r$. Тогда при любом $\beta > 0$ для


всех $x,y \in \wh D_p^n$ справедливо неравенство


\begin{equation}


\|S_d^0 x - S_d^0 y\|_{B^1} \le \exp \{\beta \delta \}\|x - y\|_


{\wh D_{2p}^n},\quad d = 1,\dotsc,r.


\end{equation}


\end{lem}





{\bf Доказательство.}


В силу леммы 3 для любых


$x \in \wh D_{2p}^n$, $d$ ($1 \le d \le r$) функция


$S_d^0 x$ измерима. Более того,


очевидно, что функция $S_d^0 x$


принадлежит пространству $B^1$ при любом $\beta > 0$.


Пусть $x,y \in \wh D_p^n $, тогда $x(t) =


\zeta _\beta (t)x_\beta (t)$,


$y(t) = \zeta _\beta (t)y_\beta (t)$, где $\zeta _\beta (t)\Def


\exp\{-\beta t\}$, $x_\beta, y_\beta \in M_{2p}$, $x_\beta =


\col({x_1}_\beta ,\dotsc,{x_n}_\beta)$,


$y_\beta =\col({y_1}_\beta ,\dotsc,{y_n}_\beta)$,


$d_i(t)\Def |(S_d^0x_i)(t) - (S_d^0 y_i)(t)|\exp \{\beta t\}$.


Таким образом, для функции $d_i$ при всех $t \ge 0$ имеем


\begin{align*}


d_i(t) &\le \exp \{\beta t\}|(S_d^0\zeta _\beta {x_i}_\beta )(t) -


(S_d^0 \zeta _\beta {y_i}_\beta )(t)| \le


\exp \{\beta t\}(S_d^0\zeta _\beta |{x_i}_\beta - {y_i}_\beta |)(t) \le \\


%


&\le\exp \{\beta t\}(S_d^0\zeta _\beta )(t)(S_d^0


|{x_i}_\beta - {y_i}_\beta |)(t) \le


\exp \{\beta \delta \}(S_d^0


|{x_i}_\beta - {y_i}_\beta |)(t)


\end{align*}


при $i = 1,\dotsc,n$, $d = 1, \dotsc, r$.


Тогда при $d = 1, \dotsc, r$


\begin{align*}


\|S_d^0 x - S_d^0 y\|_{B^1} &\le


\underset{t\ge 0}{\text{vrai \, sup}}{}


{(E|\col(d_1(t),\dotsc,d_n(t))|^{2p})}^{1/2p} \le \\


%


&\le \exp \{\beta \delta \}\sup_{t \ge 0}


(E|x_\beta (t) - y_\beta (t)|^{2p})^{1/2p} =


\exp \{\beta \delta \}\|x - y\|_{\wh D_p^n}.\quad\square


\end{align*}


\medskip





Неравенство (12) позволяет применить теорему $1'$ к


уравнению (11). Заметим, что


если уравнение (3)


$D_{2p}^n$-устойчиво, то оператор Коши $K$ этого уравнения


действует из пространства $B$ в пространство $D_{2p}^n$ и ограничен.





Пусть в дальнейшем оператор $T$ действует из пространства


$D_p^n\times (B^1)^r$ в пространство $B$. Тогда из теоремы


$1'$ получим


\medskip





{\bf Следствие.} Пусть при некотором $\beta \in [0, \ov\alpha[$ уравнение


(3) $D_{2p}^n$-устойчиво,


$\wh c\Def\|K\|_{B\to D_{2p}^n}$. Кроме того, $T\{0,\dotsc,0\}= 0$


и существуют такие постоянные $\delta'>0$,


$\delta > 0$, $k_i > 0$ $(i = 0,\dotsc,r)$, что $t - h_{id}(t) \le \delta $


при всех $t \ge 0$ и всех $i = 1,\dotsc,n$, $d = 1,\dotsc,r$,


$$


\|T\{u_0^2,\dotsc,u_r^2\} - T\{u_0^1,\dotsc,u_r^1\}\|_B \le


k_0\|u_0^2 - u_0^1\|_{D_{2p}^n} + \sum_{d=1}^r


k_d\|u_d^2 - u_d^1\|_{B^1}


$$


для всех $u_0^i \in D_{2p}^n$,


$\|u_0^i\|_{D_{2p}^n} \le \delta'$ для $i = 1, 2$,


$u_d^i \in B^1$, $\|u_d^i\|_{B^1} \le \delta'$


для $i = 1, 2$, $d = 1,\dotsc,r$, причем


$\wh c \Big(k_0 + \exp \{\beta \delta \} \sum \limits _{d=1}^rk_d\Big)<1$.


Тогда при данном $\beta$ уравнение


$dx(t)=[(Vx)(t)+(T\{x,S_1^0x,\dotsc,S_r^0x\})(t) + f(t)]dZ(t)$, $t \ge 0$,


локально $D_{2p}^n$-устойчиво.


\medskip





В качестве примера рассмотрим скалярное уравнение (11), в котором


$(Vx)(t){=}(-qx(t),0,...,0)$, $(Fx)(t)=\Big(\sum \limits _{d=1}^{r_1}


q_{1d}((S_{1d}^0x)(t)),\dotsc, \sum \limits_{d=1}^{r_m}q_{md}(


(S_{md}^0x)(t ))\Big)$,


$(S_{jd}^0x)(t)\Def


\sup \limits _{s\in [h_{jd}(t),g_{jd}(t)]}{(E|\wh x(s)|^{2p})}^{1/2p}$,


$$


\wh x(s) =


\begin{cases}


x(s), &\text{если \;} s \ge 0;\\


0, & \text{если \;} s < 0,


\end {cases}


$$


$h_{jd},g_{jd}:[0, +\infty [ \to R^1$ --- измеримые


функции, $h_{jd}(t) \le g_{jd}(t)$ и $t-h_{jd}(t) \le \delta$


при всех $t \ge 0$ и некотором $\delta > 0$, $q_{jd}$ --- некоторая


функция c $q_{jd}(0) = 0$ при $j = 1,\dotsc,m$,


$d = 1,\dotsc,r_j$, $q$ --- некоторое положительное число.





В качестве (1) возьмем уравнение (10) с


$\xi (t) = 1$, $t \ge 0$, $B = (\Lambda _{2p, \infty }^1)^\gamma$,


$\gamma (t) = \exp \{\beta t\}$, $\beta \in [0, \ov \alpha [$. Тогда


$D_{2p}^1 \subset M_{2p}^\gamma$, уравнение


(3) $D_{2p}^1$-устойчиво


при некотором $\beta \in [0, q[$ $(q < \ov \alpha )$ и


$\wh c\Def\|K\|_{B\to D_{2p}^1} = 1/(q - \beta ) +


c_p(m - 1)/(2(q - \beta))^{1/2}$.


Кроме того, если для всех $u^2, u^1 \in B$ имеем


$|q_{jd}(u^2) - q_{jd}(u^1)| \le \wh q_{jd}|u^2 - u^1|$ при всех


$j = 1,\dotsc,m$, $d = 1,\dotsc,r_j$ и выполнено неравенство


$$


q/(1 + c_p(m - 1)\sqrt {q/2}) > \sum_{d=1}^{r_1}\wh q_{1d} +


c_p\sum \limits _{j=2}^m \sum_{d=1}^{r_m}\wh q_{jd}^2\Def A,


$$


то при достаточно малом $\beta > 0$ выполняется


неравенство $\wh с \exp \{\beta \delta \}A<1$,


которое гарантирует в силу следствия при некотором


$\beta > 0$ локальную $D_{2p}^1$-устойчивость уравнения (11).





\begin{thebibliography}{9}





\bibitem{1}


Азбелев Н.В., Ермолаев М.Б., Симонов П.М. {\em К вопросу об


устойчивости


функционально-дифференциальных уравнений по первому приближению}


// Изв. вузов. Математика. -- 1995. -- \No\,10. -- С.\,3--9.


\bibitem{2}


Липцер Р.Ш., Ширяев А.Н. {\em Теория мартингалов}. -- М.:


Наука, 1986. -- 512~с.


\bibitem{3}


Jacod J. {\em Integrales stochastiques par rapport a une cemimartingale


vektorielle et changements de filtration} // Lect. Notes Math. --


1980. -- V.\,784. -- P.\,161--172.


\bibitem{4}


Кадиев Р.И., Поносов А.В. {\em Устойчивость линейных


стохастических функционально-дифференциальных


уравнений при постоянно действующих возмущениях}


// Дифференц. уравнения. -- 1992. -- Т.\,28. -- \No\,2. -- С.\,198--207.


\bibitem{5}


Кадиев Р.И. {\em Достаточные условия устойчивости стохастических систем}


// Дифференц. уравнения. -- 1994. -- Т.\,30. -- \No\,4. -- С.\,555--564.


\bibitem{6}


Хатсон В., Пим Дж.С. {\em Приложения функционального анализа


и теории операторов}. -- М.: Мир, 1983. -- 432~с.





\end{thebibliography}





\vskip 2cm





\post{\it Дагестанский государственный}{\it Поступила}


\post{\it университет}{29.10.1997}





\label{end}


\end{document}


